% SCOPE: 5 Discussion: limitations argued as claims about validity (SFT
% capability cost read against the twin, not no-SFT baselines; verify-SFT
% boilerplate bounded by the placebo arm; token-budget confound; scope and
% measurement limits; threats to validity; what a clean null still teaches);
% 6 Conclusion (one paragraph). Pre-results draft: no fabricated numbers.
\section{Discussion}
\label{sec:discussion}

\paragraph{Read absolute accuracy against the twin, not against no-SFT baselines.}
Absolute accuracy in this paper should be judged against a matched twin, not
against reinforcement-learning runs that skip supervised priming. The reason is
a known cost. Our protocol installs the meta convention through two supervised
fine-tuning stages before reinforcement learning begins. Supervised priming of
this kind lowers raw capability relative to reinforcement learning applied
directly to the base model. A public number makes the gap concrete. On
Qwen3-8B-Base with no supervised priming, \citet{kim2026rebellious} report a
standard reinforcement-learning run at 83.6 on MATH500 and 19.8 on AIME24, and
their skepticism-trained variant at 84.4 and 27.9, under a 20{,}480-token budget
averaged over sixteen samples.\footnote{These four numbers are quoted from the
cited preprint and have not been independently reproduced by us.} Our setup
starts from a post-trained checkpoint, adds two supervised stages, and trains at
a shorter budget, so it is not directly comparable to those figures. The claim
we defend is not a record score. It is the causal contribution of the meta
mechanism, measured against a twin that shares the data, the hyperparameters,
and the same priming cost. A raw-model reference row and a priming-only row
expose that cost rather than hide it.

\paragraph{The verify-SFT corpus carries boilerplate that reinforcement learning can amplify.}
The supervised corpus that teaches verification contains formulaic phrasing.
Many verification blocks restate the previous step instead of checking it
independently. Such phrasing rewards the appearance of verification over the
act. Reinforcement learning on top of this prime can enlarge the template rather
than the behavior, even when the reward scores genuine movement of evidence
\citep{gandhi2025cognitive}. A revised corpus is planned, built from traces
where verification visibly changed the outcome. The pending placebo arm bounds
how much of any observed gain the boilerplate alone can explain.

\paragraph{Token budget is a genuine confound.}
Meta blocks cost tokens, so any comparison between a meta model and a base model
mixes reasoning quality with inference-time compute. Some legacy evaluations ran
under a short token budget that cut meta traces off more often than plain ones.
Those results are superseded by a longer fixed budget. A budget-control column
then lets the reader separate thinking better from thinking longer.

\paragraph{Scope and measurement bound the claim.}
Every experiment uses one model family at one scale and one problem domain,
competition-style mathematics. Whether the reward transfers across scales,
across families, or to domains without an exact checker is untested. Grading
adds a second dependence. An earlier grader wrongly rejected many correct
base-model answers, and although the pipeline was audited afterward, residual
grading asymmetries between meta-bearing and plain traces cannot be fully ruled
out. The reward signal adds a third dependence. It rests on a frozen reference
model and on rule-based decoys. A decoy family that is too easy, or that
correlates with the model's own answer, would let the reward be satisfied
without real self-verification. That risk is why we report the SAVE and DERAIL
flip rates next to accuracy.

\paragraph{Threats to validity.}
The twin design controls everything we can audit byte for byte, but it cannot
control differences that the reward itself induces during exploration, such as
diverging entropy trajectories. We therefore compare checkpoints at the same
training step and the same seed rather than picking a best checkpoint per arm. A
second threat is statistical. Meta emission and its payoff come apart: nearly
every response carries a meta block regardless of difficulty, while the accuracy
benefit concentrates on the hardest problems. Comparing meta-bearing against
meta-free responses inside one model therefore conditions on an endogenous emit
decision and is confounded. For this reason every headline comparison is drawn
between arms and stratified by difficulty. Uniform emission on easy problems,
where a verification detour cannot help, is itself an efficiency cost we flag for
future selective-emission work.

\paragraph{A clean null would still establish a mechanism.}
Suppose the main comparison shows no accuracy gain over the matched base. The
analysis suite still answers questions the field lacks instruments for. It shows
whether explicit meta blocks move the model's own gold-versus-decoy evidence at
all. It shows whether any effect is driven by content or by form. It shows
whether the supervised prime or the reinforcement reward is the inert part. It
shows whether apparent meta effects in aggregate are artifacts of
stratification. A clean null with this apparatus would be a mechanism paper. It
would give evidence that reward routed precisely onto metacognitive content
still fails to install functional self-verification, sharpening rather than
merely echoing an existing skeptical result \citep{huang2024llmcannot}.

\section{Conclusion}
\label{sec:conclusion}
We gave an operational definition of in-reasoning metacognition. It is an
explicit \metaopen{}\dots\metaclose{} block in which the model decides whether to
verify a step or redirect its approach. We paired it with \pmishift{}, a reward
that scores each block by how far it moves the model's own evidence away from a
decoy and toward the gold answer, and that is applied to meta content alone. We
evaluated the reward against a byte-audited twin trained from the same base,
which isolates what the mechanism causes. The four-question program, covering the
main effect, the mechanism decomposition, difficulty and out-of-distribution
stratification, and calibration, is informative on either outcome: a positive
result localizes where metacognition pays, and a negative one attributes the
failure to priming, reward, or content. Results are pending
(\todo{finalize once T1--T4 land}). The next steps are to isolate the
\pmishift{} signal with a revised functional-verify corpus and additional seeds,
and to extend beyond mathematics to domains where a learned verifier, not an
exact grader, must supply the notion of gold.
